\documentclass[12pt,a4paper]{ctexbook}
\usepackage{geometry}
\geometry{margin=2.2cm}
\usepackage{setspace}
\setstretch{1.35}
\usepackage{hyperref}
\title{全唐诗选·poet.tang.3000}
\author{古典文献整理}
\date{}
\begin{document}
\maketitle
\tableofcontents
\newpage
\section{南陽道中作}
\textbf{張九齡}\par\vspace{0.5em}
登郢屬歲陰，及宛懵所適。\par
復聞東漢主，遺此南都迹。\par
佳氣藹厥初，霸圖紛在昔。\par
茲邦稱貴近，與世嘗薰赫。\par
遭遇感風雲，變衰空草澤。\par
不識鄧公樹，猶傳陰后石。\par
驅馬歷闉闍，荆榛翳阡陌。\par
事去物無象，感來心不懌。\par
懷古對窮秋，興言傷遠客。\par
眇默遵岐路，辛勤弊行役。\par
雲雁號相呼，林麕走自索。\par
顧憶徇書劍，未嘗安枕席。\par
豈暇墨突黔，空持遼豕白。\par
迷復期非遠，歸歟賞農隙。\par

\section{湘中作}
\textbf{張九齡}\par\vspace{0.5em}
湘流繞南嶽，絕目轉青青。\par
懷祿未能已，瞻途屢所經。\par
煙嶼宜春望，林猿莫夜聽。\par
永路日多緒，孤舟天復冥。\par
浮沒從此去，嗟嗟勞我形。\par

\section{彭蠡湖上}
\textbf{張九齡}\par\vspace{0.5em}
沿涉經大湖，湖流多行泆。\par
決晨趨北渚，逗浦已西日。\par
所適雖淹曠，中流且閑逸。\par
瑰詭良復多，感見乃非一。\par
廬山直陽滸，孤石當陰術。\par
一水雲際飛，數峰湖心出。\par
象類何交乣，形言豈深悉。\par
且知皆自然，高下無相恤。\par

\section{入廬山仰望瀑布水}
\textbf{張九齡}\par\vspace{0.5em}
絕頂有懸泉，喧喧出煙杪。\par
不知幾時歲，但見無昏曉。\par
閃閃青崖落，鮮鮮白日皎。\par
灑流濕行雲，濺沫驚飛鳥。\par
雷吼何噴薄，箭馳入窈窕。\par
昔聞山下蒙，今乃林巒表。\par
物情有詭激，坤元曷紛矯。\par
默然置此去，變化誰能了。\par

\section{出爲豫章郡途次廬山東巖下}
\textbf{張九齡}\par\vspace{0.5em}
茲山鎮何所，乃在澄湖陰。\par
下有蛟螭伏，上與虹蜺尋。\par
靈仙未始曠，窟宅何其深。\par
雙闕出雲峙，三宮入煙沈。\par
攀崖猶昔境，種杏非舊林。\par
想像終古跡，惆悵獨往心。\par
紛吾嬰世網，數載忝朝簪。\par
孤根自靡託，量力況不任。\par
多謝周身防，常恐橫議侵。\par
豈匪鵷鴻列，惕如泉壑臨。\par
迨茲刺江郡，來此滌塵襟。\par
有趣逢樵客，忘懷狎野禽。\par
棲閑義未果，用拙歡在今。\par
願言答休命，歸事丘中琴。\par

\section{巡屬縣道中作}
\textbf{張九齡}\par\vspace{0.5em}
春令夙所奉，駕言遵此行。\par
途中却郡掾，林下招村氓。\par
至邑無紛劇，來人但歡迎。\par
豈伊念邦政，爾實在時清。\par
短才濫符竹，弱歲起柴荆。\par
再入江村道，永懷山藪情。\par
矧逢陽節獻，默聽時禽鳴。\par
迹與素心別，感從幽思盈。\par
流芳日不待，夙志蹇無成。\par
知命且何欲，所圖唯退耕。\par
華簪極身泰，衰鬢慙木榮。\par
苟得不可遂，吾其謝世嬰。\par

\section{夏日奉使南海在道中作}
\textbf{張九齡}\par\vspace{0.5em}
緬然萬里露，赫曦三伏時。\par
飛走逃深林，流爍恐生疵。\par
行李豈無苦，而我方自怡。\par
肅事誠在公，拜慶遂及私。\par
展力慙淺效，銜恩感深慈。\par
且欲湯火蹈，況無鬼神欺。\par
朝發高山阿，夕濟長江湄。\par
秋瘴寧我毒，夏水胡不夷。\par
信知道存者，但問心所之。\par
呂梁有出入，乃覺非虛詞。\par

\section{巡按自灕水南行}
\textbf{張九齡}\par\vspace{0.5em}
理棹雖云遠，飲冰寧有惜。\par
況乃佳山川，怡然傲潭石。\par
奇峰岌前轉，茂樹隈中積。\par
猿鳥聲自呼，風泉氣相激。\par
目因詭容逆，心與清暉滌。\par
紛吾謬執簡，行郡將移檄。\par
即事聊獨歡，素懷豈兼適。\par
悠悠詠靡盬，庶以窮日夕。\par

\section{使還都湘東作}
\textbf{張九齡}\par\vspace{0.5em}
倉庚昨歸候，陽鳥今去時。\par
感物遽如此，勞生安可思。\par
養真無上格，圖進豈前期。\par
清節往來苦，壯容離別衰。\par
盛明非不遇，弱操自云私。\par
孤楫清川泊，征衣寒露滋。\par
風朝津樹落，日夕嶺猿悲。\par
牽役而無悔，坐愁秪自怡。\par
當須報恩已，終爾謝塵緇。\par

\section{冬中至玉泉山寺屬窮陰冰閉崖谷無色及仲春行縣復往焉故有此作}
\textbf{張九齡}\par\vspace{0.5em}
靈境信幽絕，芳時重暄妍。\par
再來及茲勝，一遇非無緣。\par
萬木柔可結，千花敷欲然。\par
松間鳴好鳥，竹下流清泉。\par
石壁開精舍，金光照法筵。\par
真空本自寂，假有聊相宣。\par
復此灰心者，仍追巢頂禪。\par
簡書雖有畏，身世亦相捐。\par

\section{郢城西北有大古冢數十觀其封域多是楚時諸王而年代久遠不復可識唯直西有樊妃冢因後人爲植松柏故行路盡知之}
\textbf{張九齡}\par\vspace{0.5em}
蘋藻生南澗，蕙蘭秀中林。\par
嘉名有所在，芳氣無幽深。\par
楚子初逞志，樊妃嘗獻箴。\par
能令更擇士，非直罷從禽。\par
舊國皆湮滅，先王亦莫尋。\par
唯傳賢媛隴，猶結後人心。\par
牢落山川意，蕭疎松柏陰。\par
破牆時直上，荒徑或斜侵。\par
惠問終不絕，風流獨至今。\par
千春思窈窕，黃鳥復哀音。\par

\section{荆州作二首 一}
\textbf{張九齡}\par\vspace{0.5em}
先達志其大，求意不約文。\par
士伸在知己，已況仕於君。\par
微誠夙所尚，細故不足云。\par
時來忽易失，事往良難分。\par
顧念凡近姿，焉欲殊常勳。\par
亦以行則是，豈必素有聞。\par
千慮且猶失，萬緒何其紛。\par
進士苟非黨，免相安得羣。\par
衆口金可鑠，孤心絲共棼。\par
意忠仗朋信，語勇同敗軍。\par
古劍徒有氣，幽蘭秪自薰。\par
高秩向所忝，於義如浮雲。\par

\section{荆州作二首 二}
\textbf{張九齡}\par\vspace{0.5em}
千載一遭遇，往賢所至難。\par
問余奚爲者，無階忽上摶。\par
明聖不世出，翼亮非苟安。\par
崇高自有配，孤陋何足干。\par
遇恩一時來，竊位三歲寒。\par
誰謂誠不盡，知窮力亦殫。\par
雖至負乘寇，初無挾術鑽。\par
浩蕩出江湖，翻覆如波瀾。\par
心傷不材樹，自念獨飛翰。\par
徇義在匹夫，報恩猶一餐。\par
況乃山海澤，效無毫髮端。\par
內訟已慙沮，積毀今摧殘。\par
胡爲復惕息，傷鳥畏虛彈。\par

\section{在郡秋懷二首 一}
\textbf{張九齡}\par\vspace{0.5em}
秋風入前林，蕭瑟鳴高枝。\par
寂寞遊子思，寤歎何人知。\par
宦成名不立，志存歲已馳。\par
五十而無聞，古人深所疵。\par
平生去外飾，直道如不羈。\par
未得操割效，忽復寒暑移。\par
物情自古然，身退毀亦隨。\par
悠悠滄江渚，望望白雲涯。\par
露下霜且降，澤中草離披。\par
蘭艾若不分，安用馨香爲。\par

\section{在郡秋懷二首 二}
\textbf{張九齡}\par\vspace{0.5em}
庭蕪生白露，歲候感遐心。\par
策蹇慙遠途，巢枝思故林。\par
小人恐致寇，終日如臨深。\par
魚鳥好自逸，池籠安所欽。\par
挂冠東都門，採蕨南山岑。\par
議道誠愧昔，覽分還愜今。\par
憮然憂成老，空爾白頭吟。\par

\section{忝官二十年盡在內職及爲郡嘗積戀因賦詩焉}
\textbf{張九齡}\par\vspace{0.5em}
江流去朝宗，晝夜茲不舍。\par
仲尼在川上，子牟存闕下。\par
聖達有由然，孰是無心者。\par
一郡苟能化，百城豈云寡。\par
愛禮誰爲羊，戀主吾猶馬。\par
感初時不載，思奮翼無假。\par
閑宇常自閉，沈心何用寫。\par
攬衣步前庭，登陴臨曠野。\par
白水生迢遞，清風寄瀟灑。\par
願言采芳澤，終朝不盈把。\par

\section{二弟宰邑南海見羣鴈南飛因成詠以寄}
\textbf{張九齡}\par\vspace{0.5em}
鴻雁自北來，嗷嗷度煙景。\par
常懷稻粱惠，豈憚江山永。\par
小大每相從，羽毛當自整。\par
雙鳧侶晨泛，獨鶴參宵警。\par
爲我更南飛，因書至梅嶺。\par

\section{將發還鄉示諸弟}
\textbf{張九齡}\par\vspace{0.5em}
歲陽亦頹止，林意日蕭摵。\par
云胡當此時，緬邁復爲客。\par
至愛孰能捨，名義來相迫。\par
負德良不貲，輸誠靡所惜。\par
一木逢廈構，纖塵願山益。\par
無力主君恩，寧利客卿璧。\par
去去榮歸養，憮然歎行役。\par

\section{敘懷二首 一}
\textbf{張九齡}\par\vspace{0.5em}
弱歲讀羣史，抗迹追古人。\par
被褐有懷玉，佩印從負薪。\par
志合豈兄弟，道行無賤貧。\par
孤根亦何賴，感激此爲鄰。\par

\section{敘懷二首 二}
\textbf{張九齡}\par\vspace{0.5em}
晚節從卑秩，岐路良非一。\par
既聞持兩端，復見挾三術。\par
木瓜誠有報，玉楮論無實。\par
已矣直躬者，平生壯圖失。\par
去去勿重陳，歸來茹芝朮。\par

\section{題畫山水障}
\textbf{張九齡}\par\vspace{0.5em}
心累猶不盡，果爲物外牽。\par
偶因耳目好，復假丹青妍。\par
嘗抱野間意，而迫區中緣。\par
塵事固已矣，秉意終不遷。\par
良工適我願，妙墨揮巖泉。\par
變化合羣有，高深侔自然。\par
置陳北堂上，倣像南山前。\par
靜無戶庭出，行已茲地偏。\par
萱草憂可樹，合歡忿益蠲。\par
所因本微物，況乃憑幽筌。\par
言象會自泯，意色聊自宣。\par
對玩有佳趣，使我心渺綿。\par

\section{奉和聖製瑞雪篇}
\textbf{張九齡}\par\vspace{0.5em}
萬年春，三朝日，上御明臺旅庭實。\par
初瑞雪兮霏微，俄同雲兮蒙密。\par
此時騷切陰風生，先過金殿有餘清。\par
信宿嬋娟飛雪度，能使玉人俱掩嫮。\par
皓皓樓前月初白，紛紛陌上塵皆素。\par
昨訝驕陽積數旬，始知和氣待迎新。\par
匪惟在人利，曾是扶天意。\par
天意豈云遙，雪下不崇朝。\par
皇情玩無斁，雪委方盈尺。\par
草樹紛早榮，京坻宛先積。\par
君恩誠謂何，歲稔復人和。\par
預數斯箱慶，應如此雪多。\par
朝冕旒兮載悅，想籉笠兮農節。\par
倚瑤琴兮或歌，續薰風兮瑞雪。\par
福浸昌，應尤盛，瑞雪年年常感聖。\par
願以柏梁作，長爲柳花詠。\par

\section{奉和聖製溫泉歌}
\textbf{張九齡}\par\vspace{0.5em}
有時神物待聖人，去後湯還冷，來時樹亦春。\par
今茲十月自東歸，羽斾逶迤上翠微。\par
溫谷葱葱佳氣色，離宮奕奕叶光輝。\par
臨渭川，近天邑。\par
浴日溫泉復在茲，羣仙洞府那相及。\par
吾君利物心，玄澤浸蒼黔。\par
漸漬神湯無疾苦，薰歌一曲感人深。\par

\section{南郊太尉酌獻武舞作凱安之樂}
\textbf{張九齡}\par\vspace{0.5em}
馨香惟后德，明命光天保。\par
肅祀崇聖靈，陳信表黃道。\par
玉戚初蹈厲，金匏既靜好。\par
介福何穰穰，精誠格穹昊。\par

\section{奉和聖製經孔子舊宅}
\textbf{張九齡}\par\vspace{0.5em}
孔門太山下，不見登封時。\par
徒有先王法，今爲明主思。\par
恩加萬乘幸，禮致一牢祠。\par
舊宅千年外，光華空在茲。\par

\section{奉和聖製次瓊嶽韻}
\textbf{張九齡}\par\vspace{0.5em}
山祗亦望幸，雲雨見靈心。\par
嶽館逢朝霽，關門解宿陰。\par
咸京天上近，清渭日邊臨。\par
我武因冬狩，何言是即禽。\par

\section{奉和聖製送李尚書入蜀}
\textbf{張九齡}\par\vspace{0.5em}
眷言感忠義，何有間山川。\par
徇節今如此，離情空復然。\par
皇心在勤恤，德澤委昭宣。\par
周月成功後，明年或勞還。\par

\section{奉和聖製初出洛城}
\textbf{張九齡}\par\vspace{0.5em}
東土淹龍駕，西人望翠華。\par
山川祗詢物，宮觀豈爲家。\par
十月回星斗，千官捧日車。\par
洛陽無怨思，巡幸更非賒。\par

\section{奉和聖製途次陝州作}
\textbf{張九齡}\par\vspace{0.5em}
馳道當河陝，陳詩問國風。\par
川原三晉別，襟帶兩京同。\par
後殿函關盡，前旌闕塞通。\par
行看洛陽陌，光景麗天中。\par

\section{勅賜寧王池宴}
\textbf{張九齡}\par\vspace{0.5em}
賢王有池館，明主賜春遊。\par
淑氣林間發，恩光水上浮。\par
徒慙和鼎地，終謝巨川舟。\par
皇澤空如此，輕生莫可酬。\par

\section{天津橋東旬宴得歌字韻}
\textbf{張九齡}\par\vspace{0.5em}
清洛象天河，東流形勝多。\par
朝來逢宴喜，春盡却妍和。\par
泉鮪歡時躍，林鶯醉裏歌。\par

\section{上陽水窗旬宴得移字韻}
\textbf{張九齡}\par\vspace{0.5em}
河漢非應到，汀洲忽在斯。\par
仍逢帝樂下，如逐海槎窺。\par
春賞時將換，皇恩歲不移。\par
今朝遊宴所，莫比天泉池。\par

\section{折楊柳}
\textbf{張九齡}\par\vspace{0.5em}
纖纖折楊柳，持此寄情人。\par
一枝何足貴，憐是故園春。\par
遲景那能久，芳菲不及新。\par
更愁征戍客，容鬢老邊塵。\par

\section{剪綵}
\textbf{張九齡}\par\vspace{0.5em}
姹女矜容色，爲花不讓春。\par
既爭芳意早，誰待物華真。\par
葉作參差發，枝從點綴新。\par
自然無限態，長在豔陽晨。\par

\section{和崔黃門寓直夜聽蟬之作}
\textbf{張九齡}\par\vspace{0.5em}
蟬嘶玉樹枝，向夕惠風吹。\par
幸入連宵聽，應緣飲露知。\par
思深秋欲近，聲靜夜相宜。\par
不是黃金飾，清香徒爾爲。\par

\section{和姚令公從幸溫湯喜雪}
\textbf{張九齡}\par\vspace{0.5em}
萬乘飛黃馬，千金狐白裘。\par
正逢銀霰積，如向玉京遊。\par
瑞色鋪馳道，花文拂綵旒。\par
還聞吉甫頌，不共郢歌儔。\par

\section{立春日晨起對積雪}
\textbf{張九齡}\par\vspace{0.5em}
忽對林亭雪，瑤華處處開。\par
今年迎氣始，昨夜伴春回。\par
玉潤窗前竹，花繁院裏梅。\par
東郊齋祭所，應見五神來。\par

\section{三月三日申王園亭宴集}
\textbf{張九齡}\par\vspace{0.5em}
稽亭追往事，睢苑勝前聞。\par
飛閣凌芳樹，華池落綵雲。\par
藉草人留酌，銜花鳥赴羣。\par
向來同賞處，惟恨碧林曛。\par

\section{三月三日登龍山}
\textbf{張九齡}\par\vspace{0.5em}
伊川與灞津，今日祓除人。\par
豈似龍山上，還同湘水濱。\par
衰顏憂更老，淑景望非春。\par
禊飲豈吾事，聊將偶俗塵。\par

\section{和王司馬折梅寄京邑昆弟}
\textbf{張九齡}\par\vspace{0.5em}
離別念同嬉，芬榮欲共持。\par
獨攀南國樹，遙寄北風時。\par
林惜迎春早，花愁去日遲。\par
還聞折梅處，更有棣華詩。\par

\section{晚霽登王六東閣}
\textbf{張九齡}\par\vspace{0.5em}
試上江樓望，初逢山雨晴。\par
連空青嶂合，向晚白雲生。\par
彼美要殊觀，蕭條見遠情。\par
情來不可極，日暮水流清。\par

\section{蘇侍郎紫薇庭各賦一物得芍藥}
\textbf{張九齡}\par\vspace{0.5em}
仙禁生紅藥，微芳不自持。\par
幸因清切地，還遇豔陽時。\par
名見桐君籙，香聞鄭國詩。\par
孤根若可用，非直愛華滋。\par

\section{和黃門盧侍御詠竹}
\textbf{張九齡}\par\vspace{0.5em}
清切紫庭垂，葳蕤防露枝。\par
色無玄月變，聲有惠風吹。\par
高節人相重，虛心世所知。\par
鳳皇佳可食，一去一來儀。\par

\section{和韋尚書答梓州兄南亭宴集}
\textbf{張九齡}\par\vspace{0.5em}
棠棣聞餘興，烏衣有舊遊。\par
門前杜城陌，池上曲江流。\par
暇日嘗繁會，清風詠阻修。\par
始知西峙嶽，同氣此相求。\par

\section{答陳拾遺贈竹簪}
\textbf{張九齡}\par\vspace{0.5em}
與君嘗此志，因物復知心。\par
遺我龍鍾節，非無玳瑁簪。\par
幽素宜相重，雕華豈所任。\par
爲君安首飾，懷此代兼金。\par

\section{贈澧陽韋明府}
\textbf{張九齡}\par\vspace{0.5em}
君有百鍊刃，堪斷七重犀。\par
誰開太阿匣，持割武城雞。\par
竟與尚書佩，遙應天子提。\par
何時遇操宰，當使玉如泥。\par

\section{在洪州答綦毋學士}
\textbf{張九齡}\par\vspace{0.5em}
旬雨不愆期，由來自若時。\par
爾無言郡政，吾豈欲天欺。\par
常念涓塵益，惟歡草樹滋。\par
課成非所擬，人望在東菑。\par

\section{酬王六霽後書懷見示}
\textbf{張九齡}\par\vspace{0.5em}
雲雨俱行罷，江天已洞開。\par
炎氛霽後滅，邊緒望中來。\par
作驥君垂耳，爲魚我曝鰓。\par
更憐湘水賦，還是洛陽才。\par

\section{酬王六寒朝見詒}
\textbf{張九齡}\par\vspace{0.5em}
賈生流寓日，揚子寂寥時。\par
在物多相背，唯君獨見思。\par
漁爲江上曲，雪作郢中詞。\par
忽枉兼金訊，長懷伐木詩。\par

\section{林亭詠}
\textbf{張九齡}\par\vspace{0.5em}
穿築非求麗，幽閑欲寄情。\par
偶懷因壤石，真意在蓬瀛。\par
苔益山文古，池添竹氣清。\par
從茲果蕭散，無事亦無營。\par

\section{晨出郡舍林下}
\textbf{張九齡}\par\vspace{0.5em}
晨興步北林，蕭散一開襟。\par
復見林上月，娟娟猶未沈。\par
片雲自孤遠，叢篠亦清深。\par
無事由來貴，方知物外心。\par

\section{園中時蔬盡皆鋤理唯秋蘭數本委而不顧彼雖一物有足悲者遂賦二章 一}
\textbf{張九齡}\par\vspace{0.5em}
場藿已成歲，園葵亦向陽。\par
蘭時獨不偶，露節漸無芳。\par
旨異菁爲蓄，甘非蔗有漿。\par
人多利一飽，誰復惜馨香。\par

\section{園中時蔬盡皆鋤理唯秋蘭數本委而不顧彼雖一物有足悲者遂賦二章 二}
\textbf{張九齡}\par\vspace{0.5em}
幸得不鋤去，孤苗守舊根。\par
無心羨旨蓄，豈欲近名園。\par
遇賞寧充佩，爲生莫礙門。\par
幽林芳意在，非是爲人論。\par

\section{林亭寓言}
\textbf{張九齡}\par\vspace{0.5em}
林居逢歲晏，遇物使情多。\par
蘅茝不時與，芬榮奈汝何。\par
更憐籬下菊，無如松上蘿。\par
因依自有命，非是隔陽和。\par

\section{送使廣州}
\textbf{張九齡}\par\vspace{0.5em}
家在湘源住，君今海嶠行。\par
經過正中道，相送倍爲情。\par
心逐書郵去，形隨世網嬰。\par
因聲謝遠別，緣義不緣名。\par

\section{送姚評事入蜀各賦一物得卜肆}
\textbf{張九齡}\par\vspace{0.5em}
蜀嚴化已久，沈冥空所思。\par
嘗聞賣卜處，猶憶下簾時。\par
驅傳應經此，懷賢倘問之。\par
歸來說往事，歷歷偶心期。\par

\section{送竇校書見餞得雲中辨江樹}
\textbf{張九齡}\par\vspace{0.5em}
江水天連色，無涯淨野氛。\par
微明岸傍樹，凌亂渚前雲。\par
舉棹形徐轉，登艫意漸分。\par
渺茫從此去，空復惜離羣。\par

\section{餞王司馬入計同用洲字}
\textbf{張九齡}\par\vspace{0.5em}
元僚行上計，舉餞出林丘。\par
忽望題輿遠，空思解榻遊。\par
別筵鋪柳岸，征棹倚蘆洲。\par
獨歎湘江水，朝宗向北流。\par

\section{東湖臨泛餞王司馬}
\textbf{張九齡}\par\vspace{0.5em}
南土秋雖半，東湖草未黃。\par
聊乘風日好，來汎芰荷香。\par
蘭棹無勞速，菱歌不厭長。\par
忽懷京洛去，難與共清光。\par

\section{餞濟陰梁明府各探一物得荷葉}
\textbf{張九齡}\par\vspace{0.5em}
荷葉生幽渚，芳華信在茲。\par
朝朝空此地，采采欲因誰。\par
但恐星霜改，還將蒲稗衰。\par
懷君美人別，聊以贈心期。\par

\section{餞陳學士還江南同用徵字}
\textbf{張九齡}\par\vspace{0.5em}
荷蓧旋江澳，銜杯餞霸陵。\par
別前林鳥息，歸處海煙凝。\par
風土鄉情接，雲山客念憑。\par
聖朝巖穴選，應待鶴書徵。\par

\section{通化門外送別}
\textbf{張九齡}\par\vspace{0.5em}
屢別容華改，長愁意緒微。\par
義將私愛隔，情與故人歸。\par
薄宦無時賞，勞生有事機。\par
離魂今夕夢，先繞舊林飛。\par

\section{送楊道士往天台}
\textbf{張九齡}\par\vspace{0.5em}
鬼谷還成道，天台去學仙。\par
行應松子化，留與世人傳。\par
此地煙波遠，何時羽駕旋？\par
當須一把袂，城郭共依然。\par

\section{送楊府李功曹}
\textbf{張九齡}\par\vspace{0.5em}
平生屬良友，結綬望光輝。\par
何知人事拙，相與宦情非。\par
別路穿林盡，征帆際海歸。\par
居然已多意，況復兩鄉違。\par

\section{送宛句趙少府}
\textbf{張九齡}\par\vspace{0.5em}
解巾行作吏，尊酒謝離居。\par
修竹含清景，華池澹碧虛。\par
地將幽興愜，人與舊遊疎。\par
林下紛相送，多逢長者車。\par

\section{送韋城李少府}
\textbf{張九齡}\par\vspace{0.5em}
送客南昌尉，離亭西候春。\par
野花看欲盡，林鳥聽猶新。\par
別酒青門路，歸軒白馬津。\par
相知無遠近，萬里尚爲鄰。\par

\section{送蘇主簿赴偃師}
\textbf{張九齡}\par\vspace{0.5em}
我與文雄別，胡然邑吏歸。\par
賢人安下位，鷙鳥欲卑飛。\par
激節輕華冕，移官徇綵衣。\par
羨君行樂處，從此拜庭闈。\par

\section{送廣州周判官}
\textbf{張九齡}\par\vspace{0.5em}
海郡雄蠻落，津亭壯越臺。\par
城隅百雉映，水曲萬家開。\par
里樹桄榔出，時禽翡翠來。\par
觀風猶未盡，早晚使車回。\par

\section{郡江南上別孫侍御}
\textbf{張九齡}\par\vspace{0.5em}
雲嶂天涯盡，川途海縣窮。\par
何言此地僻，忽與故人同。\par
身負邦君弩，情紆御史驄。\par
王程不我駐，離思逐秋風。\par

\section{道逢北使題贈京邑親知}
\textbf{張九齡}\par\vspace{0.5em}
征驂稍靡靡，去國方遲遲。\par
路遶南登岸，情搖北上旗。\par
故人憐別日，旅雁逐歸時。\par

\section{江上使風呈裴宣州耀卿}
\textbf{張九齡}\par\vspace{0.5em}
江路與天連，風帆何淼然。\par
遙林浪出沒，孤舫鳥聯翩。\par
常愛千鈞重，深思萬事捐。\par
報恩非徇祿，還逐賈人船。\par

\section{溪行寄王震}
\textbf{張九齡}\par\vspace{0.5em}
山氣朝來爽，溪流日向清。\par
遠心何處愜，閑棹此中行。\par
叢桂林間待，羣鷗水上迎。\par
徒然適我願，幽獨爲誰情。\par

\section{自豫章南還江上作}
\textbf{張九齡}\par\vspace{0.5em}
歸去南江水，磷磷見底清。\par
轉逢空闊處，聊洗滯留情。\par
浦樹遙如待，江鷗近若迎。\par
津途別有趣，況乃濯吾纓。\par

\section{將至岳陽有懷趙二}
\textbf{張九齡}\par\vspace{0.5em}
湘岸多深林，青冥晝結陰。\par
獨無謝客賞，況復賈生心。\par
草色雖云發，天光或未臨。\par
江潭非所遇，爲爾白頭吟。\par

\section{使還湘水}
\textbf{張九齡}\par\vspace{0.5em}
歸舟宛何處，正值楚江平。\par
夕逗煙村宿，朝緣浦樹行。\par
于役已彌歲，言旋今愜情。\par
鄉郊尚千里，流目夏雲生。\par

\section{初發道中寄遠}
\textbf{張九齡}\par\vspace{0.5em}
日夜鄉山遠，秋風復此時。\par
舊聞胡馬思，今聽楚猿悲。\par
念別朝昏苦，懷歸歲月遲。\par
壯圖空不息，常恐髮如絲。\par

\section{自湘水南行}
\textbf{張九齡}\par\vspace{0.5em}
落日摧行舫，逶迤洲渚間。\par
雖云有物役，乘此更休閑。\par
暝色生前浦，清暉發近山。\par
中流澹容與，唯愛鳥飛還。\par

\section{初入湘中有喜}
\textbf{張九齡}\par\vspace{0.5em}
征鞍窮郢路，歸棹入湘流。\par
望鳥唯貪疾，聞猿亦罷愁。\par
兩邊楓作岸，數處橘爲洲。\par
却記從來意，翻疑夢裏遊。\par

\section{耒陽溪夜行}
\textbf{張九齡}\par\vspace{0.5em}
乘夕棹歸舟，緣源路轉幽。\par
月明看嶺樹，風靜聽溪流。\par
嵐氣船間入，霜華衣上浮。\par
猿聲雖此夜，不是別家愁。\par

\section{江上}
\textbf{張九齡}\par\vspace{0.5em}
長林何繚繞，遠水復悠悠。\par
盡日餘無見，爲心那不愁。\par
憶將親愛別，行爲主恩酬。\par
感激空如此，芳時屢已遒。\par

\end{document}